\section{الأزواج الأسية: لغة لضغط التقديرات}

تسمح اختبارات المشتقات بالحصول على تقديرات منفردة، لكن التطبيقات المتكررة تحتاج إلى لغة تسجل نوع الادخار بطريقة قابلة للنقل بين الأطوار. لهذا الغرض تستعمل نظرية المجاميع الأسية مفهوم \emph{الزوج الأسي}. لن ندخل هنا في الصياغات الأكثر عمومية أو في أفضل الأزواج المعروفة، بل نثبت إطارًا تشغيليًا محدودًا يكفي لفهم عمليتي \(A\) و\(B\).

\begin{definition}[الزوج الأسي في تطبيع أحادي البعد]
\label{def:exponent-pair-framework}
لتكن \(N\ge1\)، و\(\sigma>0\)، و\(y\ge N^\sigma\)، وليكن
\(I=(a,b]\subseteq[N,2N]\). إذا كان \(P\ge1\) عددًا صحيحًا و
\(0<c<\tfrac12\)، نرمز بـ
\(\mathcal F(N,P,\sigma,y,c)\) إلى فئة الدوال الحقيقية
\(f\in C^P(I)\) التي تحقق، لكل \(x\in I\) ولكل
\(0\le p\le P-1\)،
\[
\left|
 f^{(p+1)}(x)
 -\frac{d^p}{dx^p}\bigl(yx^{-\sigma}\bigr)
\right|
\le
c\left|
 \frac{d^p}{dx^p}\bigl(yx^{-\sigma}\bigr)
\right|.
\]
نقول إن \((\kappa,\lambda)\) زوج أسي أحادي البعد إذا كان
\[
(\kappa,\lambda)\in
\left\{0\le\kappa\le\tfrac12\le\lambda\le1,\ 
\kappa+\lambda<1\right\}
\cup
\left\{\left(\tfrac12,\tfrac12\right),(0,1)\right\},
\]
وكان لكل \(\sigma>0\) عدد صحيح
\(P=P(\kappa,\lambda,\sigma)\) وثابت
\(c=c(\kappa,\lambda,\sigma)<\tfrac12\) بحيث، بانتظام لكل
\(f\in\mathcal F(N,P,\sigma,y,c)\)،
\[
\sum_{n\in I\cap\mathbb Z} e(f(n))
\ll_{\kappa,\lambda,\sigma}
\left(\frac{y}{N^\sigma}\right)^\kappa N^\lambda
\qquad(y\ge N^\sigma).
\]
\resultid{ANT-DEF-18-01}
\citedresult[\cite{TrudgianYang2024}]
\end{definition}

هذا هو التطبيع المحدد في \cite[المعادلتان~(1.1)--(1.2)]{TrudgianYang2024}؛
أما العرض الكلاسيكي لطريقة الأزواج الأسية وعمليتيها فيراجع في
\cite[الفصل~3، الصفحات~21--37]{GrahamKolesnik1991}.

\begin{remark}
ثبّتنا هنا تطبيعًا واحدًا كي يكون التعريف قابلًا للتحقق. قد تستعمل مراجع أخرى
معلمة لرتبة الطور أو تنقل قوى من \(N\) إلى تلك المعلمة، ولذلك لا تقارن الصيغ
المختلفة قبل تحويلها إلى تطبيع مشترك.
\end{remark}

الزوج التافه \((0,1)\) يعيد حدًا من رتبة طول المجموع. أما الأزواج الأفضل فتعطي ادخارًا في مناطق مناسبة من المعلمات. ولا يدعي هذا الفصل تصنيف جميع الأزواج الأسية أو الوصول إلى الجبهة المثلى الحديثة.

\section{عملية \(A\): التفريق وخفض درجة الطور}

عملية \(A\) هي الصورة المنظمة لمتباينة فرق فان دير كوربوت. إذا كان
\[
S=\sum_{n=1}^{N}e(f(n)),
\]
فإن المتباينة في لمّة \ref{lem:van-der-corput-differencing} تحول \(|S|^2\) إلى متوسط لمجاميع ذات أطوار
\[
\Delta_h f(n)=f(n+h)-f(n).
\]
إذا كان \(f\) كثير حدود من الدرجة \(d\)، فإن \(\Delta_hf\) كثير حدود من الدرجة \(d-1\). وهذه هي الفائدة الجبرية الأساسية: الفرق لا يجعل المجموع أصغر تلقائيًا، لكنه يخفض تعقيد الطور على حساب التربيع والمتوسط في \(h\).

\begin{proposition}[عملية \(A\) في النسخة المحدودة]
\label{prop:a-process-limited}
لتكن \(1\le H\le N\)، ولنفترض أن لكل \(1\le h<H\) لدينا تقدير منتظم
\[
\left|\sum_{n=1}^{N-h}e\bigl(f(n+h)-f(n)\bigr)\right|
\le B_h.
\]
عندئذ
\[
|S|^2
\le
\frac{N+H-1}{H}
\left(
N+2\sum_{h=1}^{H-1}
\left(1-\frac{h}{H}\right)B_h
\right).
\]
وبخاصة، إذا كان \(B_h\le B\) لجميع \(h\)، فإن
\[
|S|
\ll
\left(\frac{N^2}{H}+NB\right)^{1/2},
\]
بثابت مطلق.
\resultid{ANT-PROP-18-01}
\provedhere
\end{proposition}

\begin{proof}
الصيغة الأولى هي نتيجة الأطوار من متباينة فرق فان دير كوربوت مع تعويض كل ارتباط مزاح بحده \(B_h\). وللصيغة الثانية نستخدم
\[
\sum_{h=1}^{H-1}\left(1-\frac{h}{H}\right)\le \frac H2
\]
و
\[
\frac{N+H-1}{H}\ll \frac NH+1\ll\frac NH
\]
لأن \(H\le N\). ينتج
\[
|S|^2\ll \frac NH(N+HB)=\frac{N^2}{H}+NB,
\]
ثم نأخذ الجذر التربيعي.
\end{proof}

\begin{remark}[ما الذي تثبته هذه القضية؟]
تثبت القضية خطوة التفريق الكمية نفسها، لكنها لا تثبت التحويل العام للأزواج الأسية بكل فروضه وتطبيعاته. الانتقال الصوري المعروف
\[
(\kappa,\lambda)
\stackrel{A}{\longmapsto}
\left(
\frac{\kappa}{2\kappa+2},
\frac{\kappa+\lambda+1}{2\kappa+2}
\right)
\]
يعرض هنا بوصفه نتيجة كلاسيكية مقتبسة من نظرية الأزواج الأسية، لأنه يتطلب إدخال نموذج مشتقي محدد ثم تحسين اختيار \(H\) وتتبع الحدود الطرفية بصورة موحدة.
راجع \cite[الفصل~3، الصفحات~21--37]{GrahamKolesnik1991}.
\end{remark}

\begin{example}[خفض درجة طور تكعيبي]
إذا كان
\[
f(n)=\alpha n^3+\beta n^2+\gamma n+\delta,
\]
فإن
\[
\Delta_hf(n)
=3\alpha h n^2+(3\alpha h^2+2\beta h)n+
\alpha h^3+\beta h^2+\gamma h.
\]
وهكذا تحوّل عملية \(A\) المجموع التكعيبي إلى عائلة من المجاميع التربيعية، يمكن تقديرها باختبار المشتقة الثانية أو بأداة أقوى عند الحاجة.
\end{example}

\section{عملية \(B\): التحويل الثنائي والطور الساكن}

تختلف عملية \(B\) بنيويًا عن عملية \(A\). فهي ليست مجرد تطبيق آخر لمتباينة الفرق، بل تعتمد على تحويل جمعي من نوع بواسون أو فوروني، ثم تحليل تكاملات متذبذبة قرب نقاط الطور الساكن. في الصورة الإرشادية، إذا كانت \(f'\) رتيبة وقابلة للعكس، فإن مجموعًا على المتغير \(n\) يتحول إلى مجموع ثنائي على قيم صحيحة \(m\) تقع في مدى \(f'\)، ويظهر الطور المزدوج
\[
f(x_m)-m x_m,
\qquad f'(x_m)=m,
\]
مع وزن تقوده الكمية \(|f''(x_m)|^{-1/2}\).

\begin{proposition}[عملية \(B\) في إطار الأزواج الأسية]
\label{prop:b-process-cited}
في التطبيع الكلاسيكي للأزواج الأسية، تنقل عملية \(B\) الزوج \((\kappa,\lambda)\) إلى
\[
B(\kappa,\lambda)
=
\left(\lambda-\frac12,\kappa+\frac12\right),
\]
تحت فروض النعومة والرتابة وعدم الانحلال اللازمة للتحويل الثنائي وتحليل الطور الساكن.
\resultid{ANT-PROP-18-02}
\citedresult[\cite{GrahamKolesnik1991}]
\end{proposition}

لا نثبت هذه القضية هنا؛ برهانها الكامل يحتاج إلى صيغة جمع ثنائية مع حدود طرفية، وإلى تقدير موحد لتكاملات الطور الساكن. حذف هذه المكونات ثم إعلان العملية نتيجة من متباينة الفرق وحدها سيكون فجوة منطقية. لذلك يبقى تصنيفها \texttt{CITED / EXPLAINED}؛ راجع \cite[الفصل~3، الصفحات~21--37]{GrahamKolesnik1991}.

\section{تركيب العمليتين وحدود الاستخدام}

بدءًا من الزوج التافه \((0,1)\)، تعطي عملية \(B\) الزوج
\[
B(0,1)=\left(\frac12,\frac12\right).
\]
ثم تعطي عملية \(A\)، في الصياغة الكلاسيكية العامة المقتبسة، الزوج
\[
A\left(\frac12,\frac12\right)
=
\left(\frac16,\frac23\right).
\]
هذا المثال يوضح كيف تتحول الأدوات التحليلية إلى حساب على أزواج من الأسس. لكنه لا يعني أن كل تطبيق يختزل إلى اختيار أصغر \(\kappa\) أو \(\lambda\) منفردًا؛ الأفضلية تعتمد على المنطقة التي تعيش فيها معلمات الطور وطول المجموع.

\begin{remark}[صلة الفصل السابع عشر]
الأزواج الأسية أداة من أدوات ضبط المجاميع على الأقواس الصغرى، لكنها ليست وحدها البرهان الكامل لتقديرات وارينغ أو غولدباخ. يلزم في التطبيقات الفعلية تقسيم المعلمات، والتقريبات الديوفانتية، وتقديرات متوسطية، وأحيانًا مدخلات أعمق. لذلك تبقى النتائج المركبة في الفصل السابع عشر على تصنيفها المقتبس، ولا تُرقّى تلقائيًا بسبب عرض عمليتي \(A\) و\(B\) هنا.
\end{remark}

\section*{حالة دفعة التأليف الثالثة}

\begin{center}
\begin{tabular}{ll}
\texttt{ANT-ID-18-01} & \texttt{AUTHORED-DRAFT / PROVED-HERE}\\
\texttt{ANT-LEM-18-01} & \texttt{AUTHORED-DRAFT / PROVED-HERE}\\
\texttt{ANT-LEM-18-02} & \texttt{AUTHORED-DRAFT / PROVED-HERE}\\
\texttt{ANT-THM-18-01} & \texttt{AUTHORED-DRAFT / PROVED-HERE}\\
\texttt{ANT-THM-18-02} & \texttt{AUTHORED-DRAFT / CITED-EXPLAINED}\\
\texttt{ANT-DEF-18-01} & \texttt{AUTHORED-DRAFT / CITED-FRAMEWORK}\\
\texttt{ANT-PROP-18-01} & \texttt{AUTHORED-DRAFT / PROVED-HERE-LIMITED}\\
\texttt{ANT-PROP-18-02} & \texttt{AUTHORED-DRAFT / CITED-EXPLAINED}
\end{tabular}
\end{center}

جميع المعرّفات الثمانية أصبحت مؤلفة في المتن، لكنها تبقى \texttt{NON-CITABLE} إلى أن ينجح البناء، وتكتمل المراجعة الرياضية والمرجعية المستقلة، ويصدر قرار اعتماد لاحق.
